\documentclass{article}
\usepackage[T1]{fontenc}
\usepackage{lmodern}
\usepackage{graphicx}
\usepackage{amsmath,amssymb}
\usepackage[a4paper,margin=2cm]{geometry}
\usepackage[absolute,overlay]{textpos}
\setlength{\TPHorizModule}{1mm}
\setlength{\TPVertModule}{1mm}
\usepackage[indonesian]{babel}
\usepackage{hyperref}
\setlength{\emergencystretch}{2em}
% unit: Halaman 34

\begin{document}

% === Halaman 34 (python/native) ===

Pembentukan Kunci Putaran

• Setiap putaran di dalam algoritma Rijndael menggunakan kunci putaran atau 

round key. Kunci putaran dibangkitkan dari kunci eksternal dari pengguna 

yang dinamakan cipher key dan disimbolkan dengan peubah key. 

Pembangkitan semua kunci putaran dilakukan oleh fungsi KeyExpansion(). 

• Pembangkitan kunci putaran di dalam Algoritma Rijndael tergolong rumit 

dan agak sukar diterangkan. Tinjau sebuah larik key yang panjangnya 16 byte 

(16 elemen) dan Nr = 10 putaran. Sepuluh kunci putaran akan disimpan di 

dalam larik rk. Elemen awal key langsung menjadi rk[0]. Elemen-elemen 

kunci lainya akan disimpan di dalam rk[1], rk[2], …, rk[10].

\end{document}
